\begin{abstract}
\authornote{Target $\le$200 words, five moves, one sentence each.
(1) \emph{Context:} training-free diffusion samplers reconstruct accelerated MRI with
one unconditional prior, but enforce data consistency with an \emph{isotropic}
Gaussian approximation of $p(x_0\mid x_t)$
\citep{song2023pigdm,chung2023dps,peng2024optimal}.
(2) \emph{Observation:} on raw multi-coil fastMRI these samplers lose quality as the
number of steps $T$ grows ($-$\xx\,dB from $T{=}20$ to $T{=}100$ for tuned \pigdm{}).
(3) \emph{Method:} because the Cartesian MRI operator is diagonal in k-space, we
replace the isotropic covariance with a per-frequency posterior variance estimated from
the scan's own calibration region. This yields a training-free, plug-in update that
costs \xx\% extra time.
(4) \emph{Results:} across \xx{} samplers, brain and knee, and real 0.3\,T low-field
data, it improves PSNR by \pxx\,dB (95\% CI \ci) at $T{=}100$, removes the step-count
degradation, and reduces error inside fastMRI+ pathology boxes by \xx\%.
(5) \emph{Take-away:} what to report for any diffusion MRI sampler (a step-count curve)
and code availability.
Do NOT write ``scaling law''; say ``step-robust''.}
\end{abstract}

\begin{keywords}
Accelerated MRI, diffusion models, posterior sampling, data consistency, k-space,
low-field MRI
\end{keywords}
